%=============================================================================!
% 12.6  EQUAL SHARES AND THE KINETIC TEMPERATURE                              !
%=============================================================================!
\section{Equal shares and the kinetic temperature}
\label{sec:sm-equipartition}

\Cref{sec:en-friction} said that the warmth of rubbing surfaces is the
hidden energy of their jiggling atoms. This section finds how that energy
is shared among the ways the atoms can move, and from the share of the
kinetic energy defines the temperature that a simulation measures.

\subsection*{The equipartition theorem}

We take $x$ to be any one coordinate or momentum of a system in the
canonical ensemble, whose range runs from $x_1$ to $x_2$, and require
$x\,\mathrm{e}^{-\beta H}$ to vanish at both ends, as it does for a
momentum, whose range is unbounded and where the exponential falls faster
than $x$ grows, or for a coordinate held in by a potential that rises
without limit, walls included. The average of $x$ times the slope of $H$
along it is
\begin{equation*}
   \Bigl\langle x\,\frac{\partial H}{\partial x}\Bigr\rangle
   = \frac1Z\int x\,\frac{\partial H}{\partial x}\,\mathrm{e}^{-\beta H}
   \,\dd\Gamma ,
\end{equation*}
and the integral over $x$, with every other coordinate and momentum held
fixed, can be done by parts, since by the chain rule
$\partial(\mathrm{e}^{-\beta H})/\partial x = -\beta\,(\partial H/\partial
x)\,\mathrm{e}^{-\beta H}$:
\begin{align*}
   \int_{x_1}^{x_2}x\,\frac{\partial H}{\partial x}\,\mathrm{e}^{-\beta H}
   \,\dd x
   &= -\frac1\beta\int_{x_1}^{x_2}x\,\frac{\partial(\mathrm{e}^{-\beta H})}
   {\partial x}\,\dd x\\
   &= -\frac1\beta\Bigl[x\,\mathrm{e}^{-\beta H}\Bigr]_{x_1}^{x_2}
   + \frac1\beta\int_{x_1}^{x_2}\mathrm{e}^{-\beta H}\,\dd x
   && \text{parts, with $f = x$,}\\
   &= \kB T\int_{x_1}^{x_2}\mathrm{e}^{-\beta H}\,\dd x
   && \text{the bracket vanishes.}
\end{align*}
Integrating over the other coordinates and momenta too, and dividing by
$Z = \int\mathrm{e}^{-\beta H}\,\dd\Gamma$,
\begin{equation}
   \Bigl\langle x\,\frac{\partial H}{\partial x}\Bigr\rangle = \kB T ,
   \label{eq:sm-equipartition}
\end{equation}
whatever $x$ is and whatever the rest of $H$. When $H$ contains $x$ only
through a term $\frac12ax^2$, then $x\,\partial H/\partial x = ax^2$, twice
that term, and the term has the average $\frac12\kB T$. Every term of the
energy that is the square of a coordinate or momentum appearing in no
other term holds, on average, $\frac12\kB T$, whatever its coefficient:
this is the equipartition of energy. A square whose variable appears
elsewhere too, as $x$ does in $\frac12x^2 + \frac12y^2 + cxy$, need not. Chapter~14 uses \cref{eq:sm-equipartition} for the positions to
find the pressure.

Three cases recur in this book.
\begin{itemize}
\item Each component of momentum enters $H$ as $p_{ix}^2/2m_i$, so, when
   the total momentum is free, $\ens{\frac12m_iv_{ix}^2} = \frac12\kB T$
   for every atom and direction,
   whatever its mass and whatever the forces on it. At \SI{300}{\kelvin}
   that is \SI{12.93}{\milli\electronvolt}, the share given to each
   freedom of the water of \cref{sec:cs-water}.
\item A vibration of small amplitude has $H = p^2/2m + \frac12kx^2$, two
   squares, and holds $\frac12\kB T$ as kinetic and $\frac12\kB T$ as
   potential energy, $\kB T$ in all, as \cref{sec:sm-canonical} found from
   its partition function. \Cref{sec:os-energy} found the same equal
   division as an average over one period of one oscillator; here it is an
   average over the ensemble.
\item A molecule of two atoms turning freely has, in the angle $\theta$
   of its axis from the $z$ axis and the angle $\phi$ about it, $H = p_\theta^2/2I + p_\phi^2/(2I\sin^2\theta)$
   (\cref{ex:sm-rotor}). Each term is the square of a momentum, so each
   holds $\frac12\kB T$, and the turning holds $\kB T$ in all, as
   \cref{sec:ro-molecules} assumed.
\end{itemize}
A term that is not a square does not hold $\frac12\kB T$: the potential
energy of a liquid, made of Lennard-Jones terms, has no simple share. The
kinetic energy is always a sum of squares, which is why it serves as the
thermometer.

\subsection*{The kinetic temperature}

With $N_{\mathrm f}$ freedoms, each holding $\frac12\kB T$ of kinetic
energy on average, $\ens{K} = \frac12N_{\mathrm f}\kB T$, and the
instantaneous kinetic temperature
\begin{equation}
   T_{\mathrm{kin}} = \frac{2K}{N_{\mathrm f}\kB}
   \label{eq:sm-kinetic-temperature}
\end{equation}
has the temperature $T$ as its average. It is what a simulation reports
as `the temperature' at each step; it fluctuates, by an amount
\cref{sec:sm-kinetic} finds. In an isolated system the kinetic
temperature's average differs from the $T$ of \cref{eq:sm-temperature} by
terms of order $1/N$: for the ideal gas of \cref{sec:sm-micro}, $\kB T =
E/(\frac32N - 1)$, while $2K/(3N\kB)$ with $K = E$ gives $\kB
T_{\mathrm{kin}} = E/\frac32N$, smaller by the factor $1 - 2/(3N)$.

Counting $N_{\mathrm f}$ is where errors enter. $N$ atoms have $3N$
components of velocity, and each independent restriction on them removes one.
\begin{itemize}
\item A periodic run whose drift has been removed keeps its total
   momentum at zero (\cref{sec:md-start}), so its velocities satisfy three
   conditions, $\sum_im_i\vb{v}_i = \vb{0}$, and the kinetic energy is
   shared among $3N - 3$ freedoms.
\item Each distance held by RATTLE removes the component of relative
   velocity along it (\cref{sec:cs-rattle}): $n_{\mathrm c}$ independent held distances
   remove $n_{\mathrm c}$ freedoms.
\item An isolated molecule whose turning has also been removed has three
   more conditions, or two if it is linear (\cref{sec:ro-freedom}),
   leaving $3N - 6$, or $3N - 5$.
\end{itemize}
In the mass-weighted components $u_k = \sqrt{m_k}\,v_k$, the kinetic
energy is $K = \frac12\sum_ku_k^2$, half the squared length of $\vb{u}$,
and each restriction is a condition that $\vb{u}$ be at right angles to a
direction fixed by the positions at that instant: for the drift along $x$, $\sum_i\sqrt{m_i}\,u_{ix} = 0$.
Writing $\vb{u} = \sum_kc_k\vb{e}_k$ along $3N$ directions of length 1 at
right angles to one another, the first $n$ chosen to span the directions
that $n$ restrictions fix, gives $K = \frac12\sum_kc_k^2$ as in
\cref{sec:os-modes}, with $c_1$ to $c_n$ held at zero: a sum of $3N - n$
squares, each holding $\frac12\kB T$. So a periodic run with its drift
removed has $N_{\mathrm f} = 3N - n_{\mathrm c} - 3$, which
\texttt{mdlab.statmech.degrees\_of\_freedom} returns: 765 for 256 argon
atoms, and 381 for the 64 rigid water molecules of Chapter~11.
\texttt{kinetic\_temperature} then applies \cref{eq:sm-kinetic-temperature}.

The count shows in the measurements of \cref{sec:sm-maxwell}. The
kinetic energy of 765 freedoms is spread over 768 components of
velocity, so the components' mean square is $765/768$ of $\kB T/m$, and
their standard deviation $\sqrt{765/768} = 0.99804$ of $\sigma_v$:
$0.99804\times0.001674 = \SI{0.001671}{\angstrom\per\femto\second}$, the
value measured. ASE's \texttt{get\_temperature} divides by $3N$ less the
freedoms its constraints remove, and does not take off three for the
drift. On the last frame of the liquid, \texttt{mdlab} gives
\SI{137.7961}{\kelvin} and ASE \SI{137.2579}{\kelvin}, in the ratio
0.996094, which is $765/768$; the two agree once they count alike.

\subsection*{Shares between species, and with the potential}

Equipartition gives every atom the same mean kinetic energy whatever its
mass, so the heavier atoms of a mixture move more slowly and are no
colder. A liquid of 256 atoms with argon's forces, half of mass
\SI{39.948}{\amu} and half of \SI{83.798}{\amu}, the mass of krypton,
started with velocities drawn at \SI{130}{\kelvin}, has its light atoms at
\SI{126}{\kelvin} and its heavy atoms at \SI{124}{\kelvin} over the last
\SI{10}{\pico\second} of \SI{20}{\pico\second}, each temperature counting
the three components of each atom of the species. Whether a gap of
\SI{2}{\kelvin} is a real difference or the noise of a short run is the
question of Chapter~15.

Equipartition also explains what \cref{sec:md-start} saw. A crystal
started on its lattice has no potential energy above its minimum, and its
vibrations, nearly harmonic, take half of the kinetic energy they are
given: each ends with $\frac12\kB T$ of each kind. The 500 argon atoms of
that section lost half of their kinetic energy within
\SI{0.1}{\pico\second}, and 256 atoms on their lattice at the spacing of
ASE's reference value\cite{Larsen2017}, $a = \SI{5.26}{\angstrom}$, a
little below the model's \SI{5.267}{\angstrom}, which squeezes the cell
slightly but leaves the lattice, by its symmetry, where every force
vanishes, given velocities at \SI{40}{\kelvin}
(a kinetic temperature of \SI{41.5}{\kelvin}, by the scatter of
\cref{sec:sm-kinetic}), settle at \SI{21.0}{\kelvin}, 0.506 of the start.
A liquid has no such simple ratio, since its potential energy is not a sum
of squares.

\begin{practice}
Velocities drawn at a temperature set only the kinetic energy. A crystal
started on its lattice ends near half the temperature asked for, and a
system started anywhere else ends somewhere that has to be found by
running. A run is at the temperature asked for only after that
temperature has been measured over the run; a thermostat (Chapter~13)
holds it there.
\end{practice}

\begin{keyidea}
Every term of the energy that is the square of a coordinate or momentum
appearing in no other term holds $\frac12\kB T$ on average. The kinetic energy is such a sum, which
makes $T_{\mathrm{kin}} = 2K/(N_{\mathrm f}\kB)$ a thermometer, with
$N_{\mathrm f} = 3N - n_{\mathrm c} - 3$ in a periodic run whose drift is
removed. Atoms of every mass share alike; a harmonic crystal shares
equally between kinetic and potential energy.
\end{keyidea}

\notebook{12}{count the freedoms of a run and compare the temperature with ASE's}

\begin{selfcheck}
A box of 100 rigid water molecules, drift removed, has a kinetic energy
of \SI{7.4}{\electronvolt}. What is its kinetic temperature?
\end{selfcheck}
